\documentclass{article}
\usepackage[margin=1in]{geometry} % margins
\usepackage[justification=centering]{caption}
%\usepackage{calrsfs}
\usepackage{amsmath}
\usepackage{amsfonts}
\usepackage{amssymb}
\usepackage{parskip}
\usepackage{graphicx}
\usepackage{enumerate}
\usepackage{enumitem}
\usepackage{amsthm}
\usepackage{float}
\usepackage{tikz-cd}
\usepackage{ytableau}
\usepackage{dsfont}

\renewcommand{\implies}{\Rightarrow}
\newcommand{\longimplies}{\Longrightarrow}
\newcommand{\N}{\mathbb{N}}
\newcommand{\Z}{\mathbb{Z}}
\newcommand{\Q}{\mathbb{Q}}
\newcommand{\R}{\mathbb{R}}
\newcommand{\C}{\mathbb{C}}
\newcommand{\D}{\mathbb{D}}
\renewcommand{\H}{\mathbb{H}}
\newcommand{\F}{\mathbb{F}}
\newcommand{\K}{\mathbb{K}}
\newcommand{\Lf}{\mathbb{L}}
\newcommand{\Sp}{\mathbb{S}}
\newcommand{\cl}[1]{\overline{#1}}
\renewcommand{\emptyset}{\varnothing}
\newcommand{\norm}[1]{\left\|#1\right\|}
\newcommand{\maxnorm}[1]{\left\|#1\right\|_{\text{max}}}
\newcommand{\opnorm}[1]{\left\|#1\right\|_{\mathrm{op}}}
\newcommand{\supnorm}[1]{\left\|#1\right\|_{\infty}}
\newcommand{\id}{\operatorname{id}}
\newcommand{\sspan}{\operatorname{span}}
\newcommand{\tr}{\operatorname{tr}}
\newcommand{\ch}{\operatorname{char}}
\newcommand{\poly}{\mathrm{P}}
\newcommand{\mtx}{\mathcal{M}}
\newcommand{\seq}{\mathbf{Seq}}
\newcommand{\vv}[1]{\mathbf{#1}}
\newcommand{\dirsum}{\oplus}
\newcommand{\rank}{\operatorname{rank}}
\newcommand{\im}{\operatorname{im}}

\newcommand{\cj}[1]{\overline{#1}}
\newcommand{\inprod}[2]{\left\langle#1,#2\right\rangle}
\renewcommand{\Re}{\operatorname{Re}}
\renewcommand{\Im}{\operatorname{Im}}

\newcommand{\res}{\operatorname{res}}

\newcommand{\floor}[1]{\left\lfloor#1\right\rfloor}

% Theorem environments
\usepackage{amsthm}

\theoremstyle{definition} % The "plain" style italicizes all body text.
	\newtheorem{thm}{Theorem}
		\numberwithin{thm}{section} % Theorem numbers are determined by section.
	\newtheorem{lemma}[thm]{Lemma}
	\newtheorem{prop}[thm]{Proposition}
	\newtheorem{cor}[thm]{Corollary}
	\newtheorem{conj}[thm]{Conjecture}

\theoremstyle{definition} % The "definition" style does not italicize body text..
    \newtheorem{defn}[thm]{Definition}
	\newtheorem{example}[thm]{Example}
    \newtheorem{exercise}[thm]{Exercise}
    \newtheorem{question}[thm]{Question}

\usepackage{mathtools}

\usetikzlibrary{decorations.markings,positioning,decorations.text}

\def\xr{5}
\def\yr{5}

%\usepackage{titlesec}
%\titleformat{\section}[hang]{\normalfont\fontsize{14}{0}\bfseries}{Problem 8.\thesection\ }{0pt}{}

\begin{document}

\section{Worksheet}

\begin{question}[Practice 1]
Let $f \in L^2$ with $\sum_n |\hat{f}(n)| < \infty$. Show that there is a continuous function $g$ such that $f = g$ a.e.
\end{question}
\begin{proof}

\end{proof}

\begin{question}[Practice 2]
Let $f : \mathbb{R} \to \mathbb{R}$ be $2\pi$-periodic and in $C^\alpha$. That is, $|f(x) - f(y)| \leq C |x - y|^\alpha$ for some $C > 0$ and $\alpha \in (0, 1)$. Show that the Fourier series of $f$ converges pointwise to $f$.
\end{question}
\begin{proof}

\end{proof}

\begin{question}[Practice 3]
Show that $L^\infty$ is complete.
\end{question}
\begin{proof}

\end{proof}

\begin{question}[Practice 4]
Let $f : E \to \mathbb{R}$. If either $m(E) < \infty$ or $f \in L^p(E)$ for some $p \geq 1$, show that
\[
\lim_{p \to \infty} \|f\|_{L^p(E)} = \|f\|_{L^\infty(E)}.
\]
Find a counterexample if the assumptions are dropped.
\end{question}
\begin{proof}

\end{proof}

\begin{question}[Practice 5]
Let $p \in [1, \infty)$. Let $A_p$ be the set of functions $f : [0, 1] \to \mathbb{R}$ that have a continuous derivative with
\[
\int_0^1 |f'(x)|^p dx \leq 1.
\]
For which $p$ are closed (with respect to the topology of uniform convergence) subsets of $A_p$ compact?
\end{question}
\begin{proof}

\end{proof}

\begin{question}[Practice 6]
Let $p \in [1, \infty]$. Construct a function $f \in L^p(\mathbb{R})$ but $f \notin L^q(\mathbb{R})$ for any $q \neq p$.
\end{question}
\begin{proof}

\end{proof}

\begin{question}[Practice 7]
Let $f_n$ be a sequence that converges in $L^p$ to some $f$. Show that there exists a subsequence that converges a.e.
\end{question}
\begin{proof}

\end{proof}

\begin{question}[Practice 8]
This problem explores the concept of \textit{convergence in measure}. Let $f_n : E \to \mathbb{R}$ be a sequence of measurable functions. We say that $f_n$ converges to some $f$ in measure if for all $\varepsilon > 0$ we have
\[
\lim_{n \to \infty} m(\{x : |f_n(x) - f(x)| > \varepsilon\}) = 0.
\]
\begin{enumerate}[label=(\alph*)]
    \item If $f_n$ converges to $f$ in $L^p$ show that it converges in measure.
    \item If a sequence converges in measure, show that a subsequence converges pointwise a.e.
    \item If $m(E) < \infty$ and $f_n$ converges pointwise, show that it converges in measure.
    \item Find sequences that converge in measure but not pointwise.
    \item Show that Fatou’s Lemma and the MCT hold if pointwise convergence is replaced by convergence in measure.
\end{enumerate}
\end{question}
\begin{proof}

\end{proof}

\begin{question}[Practice 9]
This problem shows that the Riemann integral coincides with the Lebesgue integral. In this problem let $f : [a, b] \to \mathbb{R}$ be a bounded function (not necessarily measurable).
\begin{enumerate}[label=(\alph*)]
    \item Let $P_k$ be a sequence of partitions of $[a, b]$ whose mesh goes to 0 and are nested inside each other (i.e., $P_{k+1}$ is a refinement of $P_k$ for all $k$). For $P_k$ let
    \[
    G_k(x) = \sum_j M_j \mathbf{1}_{(t_{j-1}, t_j]}, \quad g_k(x) = \sum_j m_j \mathbf{1}_{(t_{j-1}, t_j]},
    \]
    where the intervals $(t_{j-1}, t_j]$ are the partition of $[a, b]$ corresponding to $P_k$ and $M_j$, $m_j$ are the sup and inf of $f$ over $[t_{j-1}, t_j]$. Show that the Lebesgue integrals of $g_k$ and $G_k$ correspond to the upper and lower Riemann sums associated to $P_k$ and $f$.
    
    \item Let $G(x) = \lim_{k \to \infty} G_k(x)$ and $g(x) = \lim_{k \to \infty} g_k(x)$ (why do the limits exist?). Note $g \leq f \leq G$ everywhere. Show that if $f$ is Riemann integrable, then $g = G$ a.e. (simply note that the Lebesgue integrals of $G$ and $g$ are equal to the limits of the Lebesgue integrals of $G_k$ and $g_k$ which coincide if $f$ is Riemann integrable, as these integrals equal the upper and lower Riemann sums associated to the partitions $P_k$). Therefore, $f$ is measurable as $g$ and $G$ are limits of measurable functions and its Lebesgue integral equals its Riemann integral.
    
    \item Define
    \[
    H(x) = \lim_{\delta \to 0} \sup_{|y - x| \leq \delta} f(y), \quad h(x) = \lim_{\delta \to 0} \inf_{|y - x| \leq \delta} f(y).
    \]
    Why do the above limits exist? Show that $H(x) = h(x)$ iff $f$ is continuous at $x$.
    
    \item Show that the Lebesgue integrals of $H$ and $h$ are the limits of the upper and lower Riemann sums of $f$ over any sequence of partitions whose mesh goes to 0. So if $f$ is Riemann integrable, $H = h$ a.e. Therefore $f$ is the a.e. limit of continuous functions and hence is Lebesgue measurable and its Lebesgue integral equals its Riemann integral.
\end{enumerate}
\end{question}
\begin{proof}

\end{proof}
    

\end{document}